\documentclass[11pt]{article}

% Practice problems for Lecture 1
\usepackage[margin=1in]{geometry}
\usepackage[T1]{fontenc}
\usepackage{lmodern}
\usepackage{amsmath,amssymb,amsthm}
\usepackage{enumitem}
\usepackage{hyperref}
\usepackage{xcolor}

\title{Lecture 1: Practice Problems}
\author{Linear Algebra}
\date{\today}

\theoremstyle{definition}
\newtheorem{problem}{Problem}[section]

% Mandatory problems are marked in their title; change this macro to restyle all of them.
\newcommand{\mandatory}{\textcolor{red!70!black}{\textbf{Mandatory}}}

\DeclareMathOperator{\im}{im}
\DeclareMathOperator{\rank}{rank}
\DeclareMathOperator{\Null}{null}
\DeclareMathOperator{\col}{col}
\DeclareMathOperator{\row}{row}
\DeclareMathOperator{\Span}{span}
\DeclareMathOperator{\RREF}{RREF}
\DeclareMathOperator{\tr}{tr}
\DeclareMathOperator{\diag}{diag}

\begin{document}

\maketitle

\section{Number Sets and Fields}

\begin{problem}\label{pr:numbers}
    For each of the numbers
    \[\frac{7}{3}, \qquad \sqrt{49}, \qquad \sqrt{12}, \qquad 0.25, \qquad \pi - \pi, \qquad (1 + i)(1 - i), \qquad 2 + 3i,\]
    decide which of the sets $\mathbb{Z}$, $\mathbb{Q}$, $\mathbb{R}$ and $\mathbb{C}$ it belongs to.
\end{problem}

\begin{problem}[\mandatory]\label{pr:field-fail}
    Each of the following sets, with the usual addition and multiplication, is not a field. For each set, find a field axiom that fails.
    \begin{enumerate}[label=(\alph*)]
        \item $\mathbb{N} = \{0, 1, 2, \ldots\}$;
        \item the set $\mathbb{R}_{>0}$ of positive real numbers;
        \item the set $\{0, 1\}$;
        \item the set $2\mathbb{Z} = \{\ldots, -4, -2, 0, 2, 4, \ldots\}$ of even integers.
    \end{enumerate}
\end{problem}

\begin{problem}\label{pr:Qsqrt2}
    \begin{enumerate}[label=(\alph*)]
        \item Show that $\mathbb{Q}(\sqrt{2}) = \{a + b\sqrt{2} : a, b \in \mathbb{Q}\}$, with the usual addition and multiplication, is a field. (Hint: to find the inverse of $a + b\sqrt{2} \neq 0$, multiply by $\frac{a - b\sqrt{2}}{a - b\sqrt{2}}$. Why is $a^2 - 2b^2 \neq 0$?)
        \item Is $\{a + b\sqrt{2} : a, b \in \mathbb{Z}\}$ a field?
    \end{enumerate}
\end{problem}

\section{Vector Spaces}

\begin{problem}[\mandatory]\label{pr:weird-ops}
    On $\mathbb{R}^2$, keep the usual vector addition, but define a new scalar multiplication. In each case, decide which vector space axioms hold and which fail.
    \begin{enumerate}[label=(\alph*)]
        \item $c \odot (x, y) = (cx, 0)$;
        \item $c \odot (x, y) = (cx, y)$.
    \end{enumerate}
\end{problem}

\begin{problem}\label{pr:positive-reals}
    On the set $V = \mathbb{R}_{>0}$ of positive real numbers, define ``addition'' and ``scalar multiplication'' by
    \[u \oplus v = uv \qquad \text{and} \qquad c \odot u = u^c \qquad (u, v \in V, \ c \in \mathbb{R}).\]
    Show that $V$ is a vector space over $\mathbb{R}$. What is its zero vector, and what is the additive inverse of $u$?
\end{problem}

\begin{problem}\label{pr:functions}
    \begin{enumerate}[label=(\alph*)]
        \item Show that the set of all functions $f: \mathbb{R} \to \mathbb{R}$, with $(f + g)(x) = f(x) + g(x)$ and $(cf)(x) = c f(x)$, is a vector space over $\mathbb{R}$. What is its zero vector?
        \item Is the set of all functions $f: \mathbb{R} \to \mathbb{R}$ with $f(0) = 1$, with the same operations, a vector space?
    \end{enumerate}
\end{problem}

\section{Bases and Dimension}

\begin{problem}[\mandatory]\label{pr:indep}
    Decide whether the vectors $(1, 2, 1)$, $(2, 1, 0)$ and $(1, -1, -1)$ in $\mathbb{R}^3$ are linearly independent. If they are not, find a nontrivial linear combination that equals $0$, and find a basis and the dimension of their span.
\end{problem}

\begin{problem}[\mandatory]\label{pr:basisU}
    One can check that $U = \{(x, y, z, w) \in \mathbb{R}^4 : x + y + z + w = 0 \text{ and } x = w\}$ is a vector space. Find a basis of $U$ and its dimension.
\end{problem}

\begin{problem}\label{pr:taylor-basis}
    Show that $\mathcal{B} = (1, \ x - 1, \ (x - 1)^2)$ is an ordered basis of the vector space $\mathbb{R}[x]_{\leq 2}$ of polynomials of degree at most $2$. Find the coordinate vectors $[x^2]_{\mathcal{B}}$ and $[3 + 2x + x^2]_{\mathcal{B}}$.
\end{problem}

\section{Linear Transformations}

\begin{problem}[\mandatory]\label{pr:linear}
    Which of the following maps are linear transformations? Justify your answers.
    \begin{enumerate}[label=(\alph*)]
        \item $T: \mathbb{R}^2 \to \mathbb{R}^2$, $T(x, y) = (x - y, 3x)$;
        \item $T: \mathbb{R}^2 \to \mathbb{R}^2$, $T(x, y) = (x + y, 1)$;
        \item $T: \mathbb{R}^3 \to \mathbb{R}^2$, $T(x, y, z) = (xy, z)$;
        \item $D: \mathbb{R}[x]_{\leq 3} \to \mathbb{R}[x]_{\leq 3}$, $D(p) = p'$ (the derivative);
        \item $T: \mathbb{R} \to \mathbb{R}$, $T(x) = |x|$;
        \item $T: \mathbb{R}^2 \to \mathbb{R}^2$, $T(x, y) = (y, x)$.
    \end{enumerate}
\end{problem}

\begin{problem}[\mandatory]\label{pr:kerim}
    Let $T: \mathbb{R}^3 \to \mathbb{R}^3$ be defined by $T(x, y, z) = (x + y, \ y + z, \ x - z)$. Find the kernel and the image of $T$. Is $T$ injective? Is it surjective?
\end{problem}

\begin{problem}\label{pr:gallery}
    For each of the following linear transformations of the plane, compute $T(e_1)$ and $T(e_2)$, sketch the image of the unit square and of the letter F from the gallery figure in the lecture, and find the kernel and the image of $T$.
    \begin{enumerate}[label=(\alph*)]
        \item the reflection in the line $y = x$, $T(x, y) = (y, x)$;
        \item the projection onto the $x$-axis, $T(x, y) = (x, 0)$;
        \item the shear $T(x, y) = (x + 2y, y)$;
        \item $T(x, y) = (-x, -y)$.
    \end{enumerate}
    In which cases is the image of the letter F mirrored?
\end{problem}

\begin{problem}\label{pr:derivative}
    Let $D: \mathbb{R}[x]_{\leq 3} \to \mathbb{R}[x]_{\leq 3}$, $D(p) = p'$. Find the kernel and the image of $D$, and decide whether $D$ is injective or surjective. What is $D \circ D \circ D \circ D$?
\end{problem}

\begin{problem}[$\star$]\label{pr:composition}
    Let $S$, $T$, $U$ be nonempty sets, and let $f: S \to T$ and $g: T \to U$ be functions.
    \begin{enumerate}[label=(\alph*)]
        \item Show that if $g \circ f$ is injective, then $f$ is injective.
        \item Give an example in which $g \circ f$ is injective, but $g$ is not.
        \item Show that if $g \circ f$ is surjective, then $g$ is surjective.
        \item Let $A: \mathbb{R}^m \to \mathbb{R}^n$ and $B: \mathbb{R}^n \to \mathbb{R}^m$ be linear transformations with $B \circ A = \mathrm{id}_{\mathbb{R}^m}$. Which of $A$ and $B$ must be injective, and which must be surjective?
    \end{enumerate}
\end{problem}

\begin{problem}[$\star$]\label{pr:distributive}
    Let $\mathcal{F}$ be the set of all functions $\mathbb{R} \to \mathbb{R}$, with pointwise addition $(f + g)(x) = f(x) + g(x)$ and composition $(f \circ g)(x) = f(g(x))$.
    \begin{enumerate}[label=(\alph*)]
        \item Show that $\circ$ is not commutative.
        \item Show that the right distributive law $(g + h) \circ f = (g \circ f) + (h \circ f)$ always holds.
        \item Show that the left distributive law $f \circ (g + h) = (f \circ g) + (f \circ h)$ can fail. (Try $f(x) = x^2$.)
        \item Show that the left distributive law holds whenever $f$ is linear, i.e., $f(x) = cx$ for some $c \in \mathbb{R}$. Compare with the distributive laws for matrix multiplication.
    \end{enumerate}
\end{problem}

\begin{problem}[\mandatory]\label{pr:iso}
    \begin{enumerate}[label=(\alph*)]
        \item Show that the following map is an isomorphism:
        \[T: \mathbb{R}[x]_{\leq 3} \to \mathbb{R}^4, \qquad T(p) = \big(p(0), p'(0), \tfrac{1}{2}p''(0), \tfrac{1}{6}p'''(0)\big).\] (Hint: compute $T(a_0 + a_1 x + a_2 x^2 + a_3 x^3)$.)
        \item Show that the vector space $V$ of Problem~\ref{pr:positive-reals} is isomorphic to $\mathbb{R}$. (Hint: logarithm.)
        \item Which of the following pairs of real vector spaces are isomorphic: $\mathbb{R}[x]_{\leq 5}$ and $\mathbb{R}^6$; $\mathbb{C}$ (over $\mathbb{R}$) and $\mathbb{R}^3$; $\mathbb{R}^2 \times \mathbb{R}^3$ and $\mathbb{R}^5$?
    \end{enumerate}
\end{problem}

\section{Subspaces}

\begin{problem}[\mandatory]\label{pr:subspaces}
    Which of the following subsets of $\mathbb{R}^3$ are subspaces?
    \begin{enumerate}[label=(\alph*)]
        \item $\{(x, y, z) : x = 2y\}$;
        \item $\{(x, y, z) : z = x^2\}$;
        \item $\{(x, y, z) : x + y = z + 1\}$;
        \item $\{(x, y, z) : x = 0 \text{ or } y = 0\}$;
        \item $\{(t, 2t, 3t) : t \in \mathbb{R}\}$;
        \item $\{(x, y, z) : x + y + z = 0 \text{ and } x = y\}$.
    \end{enumerate}
\end{problem}

\begin{problem}\label{pr:intersection}
    Let $U$ and $W$ be subspaces of a vector space $V$.
    \begin{enumerate}[label=(\alph*)]
        \item Show that $U \cap W$ is a subspace of $V$.
        \item Show that $U \cup W$ need not be a subspace.
        \item Show that $U \cup W$ is a subspace if and only if $U \subseteq W$ or $W \subseteq U$.
    \end{enumerate}
\end{problem}

\begin{problem}\label{pr:poly-subspace}
    Show that $\{p \in \mathbb{R}[x]_{\leq 3} : p(1) = 0\}$ is a subspace of $\mathbb{R}[x]_{\leq 3}$ by writing it as the kernel of a linear transformation. Find a basis of this subspace and its dimension.
\end{problem}

\section{Matrices}

\subsection{Matrix Operations}

\begin{problem}[\mandatory]\label{pr:matvec}
    Let $A = \begin{pmatrix} 1 & 0 & 2 \\ -1 & 3 & 1 \end{pmatrix}$ and $x = (2, 1, -1)^\top$.
    \begin{enumerate}[label=(\alph*)]
        \item Compute $Ax$ in two ways: entry by entry, and as a linear combination of the columns of $A$.
        \item Compute $A^\top$, $AA^\top$ and $A^\top A$. Which of these matrices are symmetric?
    \end{enumerate}
\end{problem}

\begin{problem}\label{pr:sizes}
    Let $A$ be a $2 \times 3$ matrix, $B$ a $3 \times 3$ matrix and $C$ a $3 \times 2$ matrix. Which of the following expressions are defined, and what are their sizes?
    \[AB, \qquad BA, \qquad AC, \qquad CA, \qquad A + C^\top, \qquad A^\top + C, \qquad ABC.\]
\end{problem}

\begin{problem}\label{pr:reflect-rotate}
    Let $R = \begin{pmatrix} 0 & -1 \\ 1 & 0 \end{pmatrix}$ be the rotation by $90^\circ$, and $F = \begin{pmatrix} 1 & 0 \\ 0 & -1 \end{pmatrix}$ the reflection in the $x$-axis. Compute $FR$ and $RF$, and describe both as geometric transformations. What does this tell you about the order in which linear transformations are composed?
\end{problem}

\begin{problem}\label{pr:rotation-addition}
    Let $R_\theta = \begin{pmatrix} \cos\theta & -\sin\theta \\ \sin\theta & \cos\theta \end{pmatrix}$ be the rotation by the angle $\theta$.
    \begin{enumerate}[label=(\alph*)]
        \item Explain geometrically why $R_\alpha R_\beta = R_{\alpha + \beta}$.
        \item Multiply out $R_\alpha R_\beta$, and deduce the addition formulas for $\cos(\alpha + \beta)$ and $\sin(\alpha + \beta)$.
        \item Show that $(R_\theta)^n = R_{n\theta}$, and conclude that $(R_{2\pi/n})^n = I_2$.
    \end{enumerate}
\end{problem}

\begin{problem}[$\star$]\label{pr:complex-matrix}
    For a complex number $z = a + bi$, let $M(z) = \begin{pmatrix} a & -b \\ b & a \end{pmatrix}$.
    \begin{enumerate}[label=(\alph*)]
        \item Show that $M(z + w) = M(z) + M(w)$ and $M(zw) = M(z) M(w)$ for all $z, w \in \mathbb{C}$.
        \item Show that $M(1 + i) = \sqrt{2}\, R_{\pi/4}$, with $R_\theta$ as in Problem~\ref{pr:rotation-addition}.
        \item Use parts (a) and (b) to compute $(1 + i)^{10}$.
    \end{enumerate}
\end{problem}

\subsection{Row Reduction}

\begin{problem}\label{pr:ref}
    Which of the following matrices are in row echelon form? Which are in reduced row echelon form?
    \[\begin{pmatrix} 1 & 3 & 0 & 2 \\ 0 & 0 & 1 & 4 \\ 0 & 0 & 0 & 0 \end{pmatrix}, \quad \begin{pmatrix} 1 & 3 & 1 \\ 0 & 1 & 2 \\ 0 & 0 & 0 \end{pmatrix}, \quad \begin{pmatrix} 0 & 0 & 1 \\ 0 & 1 & 0 \\ 1 & 0 & 0 \end{pmatrix}, \quad \begin{pmatrix} 1 & 0 & 0 \\ 0 & 0 & 0 \\ 0 & 0 & 1 \end{pmatrix}, \quad \begin{pmatrix} 2 & 1 \\ 0 & 1 \end{pmatrix}.\]
\end{problem}

\begin{problem}[\mandatory]\label{pr:rref-compute}
    Find the reduced row echelon form and the rank of
    \[P = \begin{pmatrix} 1 & 2 & 1 & 0 \\ 2 & 4 & 0 & 2 \\ 1 & 2 & 3 & -2 \end{pmatrix}.\]
\end{problem}

\begin{problem}[\mandatory, $\star$]\label{pr:system}
    Consider the system of linear equations
    \begin{align*}
        w + 3x + 2y + 2z &= 0, \\
        w + 4x + y &= 0, \\
        3w + 5x + 10y + 14z &= 0, \\
        2w + 5x + 5y + 6z &= 0.
    \end{align*}
    \begin{enumerate}[label=(\alph*)]
        \item Find all solutions $(w, x, y, z) \in \mathbb{R}^4$ by bringing the coefficient matrix into reduced row echelon form.
        \item Check that $(-5, 1, 1, 0)$ and $(0, 2, -8, 5)$ are solutions. Is every solution a scalar multiple of $(-5, 1, 1, 0)$?
        \item What is the dimension of the space of solutions?
    \end{enumerate}
\end{problem}

\subsection{The Inverse of a Matrix}

\begin{problem}\label{pr:inverse2}
    Find the inverse of $A = \begin{pmatrix} 3 & 1 \\ 5 & 2 \end{pmatrix}$ by solving the equation $Ax = y$ for $x$, as in the lecture. Check your answer by computing $AA^{-1}$.
\end{problem}

\begin{problem}[\mandatory]\label{pr:inverse3}
    Use the Gauss--Jordan method to find the inverse of
    \[A = \begin{pmatrix} 1 & 2 & 1 \\ 0 & 1 & 1 \\ 1 & 1 & 1 \end{pmatrix},\]
    and use it to solve the system $Ax = (1, 2, 3)^\top$.
\end{problem}

\begin{problem}\label{pr:arithmetic}
    \begin{enumerate}[label=(\alph*)]
        \item Apply the Gauss--Jordan method to $A = \begin{pmatrix} 1 & 2 & 3 \\ 2 & 3 & 4 \\ 3 & 4 & 5 \end{pmatrix}$ to show that $A$ is singular, and find a nonzero vector $x$ with $Ax = 0$.
        \item Show more generally that every $3 \times 3$ matrix whose rows form an arithmetic progression (i.e., row $2$ $-$ row $1$ $=$ row $3$ $-$ row $2$) is singular.
    \end{enumerate}
\end{problem}

\begin{problem}\label{pr:elementary}
    Let $A = \begin{pmatrix} 0 & 2 \\ 1 & 3 \end{pmatrix}$. Find elementary matrices $E_1$, $E_2$, $E_3$ such that $E_3 E_2 E_1 A = I_2$. Use them to write $A^{-1}$, and also $A$ itself, as a product of elementary matrices.
\end{problem}

\begin{problem}[$\star$]\label{pr:idempotent}
    A square matrix $A$ with $A^2 = A$ is called idempotent.
    \begin{enumerate}[label=(\alph*)]
        \item Show that the only invertible idempotent matrix is the identity.
        \item Show that $P = \begin{pmatrix} 1 & 1 \\ 0 & 0 \end{pmatrix}$ is idempotent, and find the kernel and the image of $T_P$. Describe $T_P$ geometrically.
        \item Show that if $A$ is idempotent, then so is $I - A$, and that $A(I - A) = 0$.
        \item Let $A$ be an idempotent $n \times n$ matrix. Show that $\im(T_A) \cap \ker(T_A) = \{0\}$, and that every $v \in \mathbb{R}^n$ can be written as $v = Av + (v - Av)$, where $Av \in \im(T_A)$ and $v - Av \in \ker(T_A)$.
    \end{enumerate}
\end{problem}

\begin{problem}\label{pr:nilpotent}
    Let $N = \begin{pmatrix} 0 & 1 & 2 \\ 0 & 0 & 3 \\ 0 & 0 & 0 \end{pmatrix}$.
    \begin{enumerate}[label=(\alph*)]
        \item Compute $N^2$ and $N^3$.
        \item Show that $(I - N)(I + N + N^2) = I$, and compute $(I - N)^{-1}$.
        \item More generally, show that if $N^k = 0$ for some $k \geq 1$, then $I - N$ is invertible, with $(I - N)^{-1} = I + N + N^2 + \cdots + N^{k-1}$. (Compare with the geometric series $\frac{1}{1 - x} = 1 + x + x^2 + \cdots$.)
    \end{enumerate}
\end{problem}

\begin{problem}[\mandatory]\label{pr:interpolation}
    Let $T: \mathbb{R}[x]_{\leq 2} \to \mathbb{R}^3$, $T(p) = \big(p(-1), p(0), p(1)\big)$.
    \begin{enumerate}[label=(\alph*)]
        \item Find the matrix of $T$ with respect to the ordered basis $(1, x, x^2)$ of $\mathbb{R}[x]_{\leq 2}$ and the standard basis of $\mathbb{R}^3$.
        \item Show that this matrix is invertible, and conclude that $T$ is an isomorphism.
        \item Find the polynomial $p \in \mathbb{R}[x]_{\leq 2}$ with $p(-1) = 0$, $p(0) = 1$, $p(1) = 0$, and the polynomial $q$ with $q(-1) = 1$, $q(0) = 2$, $q(1) = 5$.
    \end{enumerate}
\end{problem}

\subsection{Special Types of Matrices}

\begin{problem}\label{pr:dims-special}
    Find the dimensions of the following vector spaces, and a basis of each of the first two:
    \begin{enumerate}[label=(\alph*)]
        \item $\mathrm{Sym}_2$;
        \item $\mathrm{Skew}_2$;
        \item $\mathbb{R}^{2 \times 3}$;
        \item $\mathrm{Sym}_4$;
        \item $\mathrm{Skew}_4$;
        \item the space of upper triangular $3 \times 3$ matrices with trace $0$.
    \end{enumerate}
\end{problem}

\begin{problem}\label{pr:groups}
    Which of the following sets of matrices are matrix groups, i.e., contain the identity and are closed under multiplication and under taking inverses?
    \begin{enumerate}[label=(\alph*)]
        \item the invertible diagonal $n \times n$ matrices;
        \item the invertible symmetric $2 \times 2$ matrices;
        \item the $2 \times 2$ matrices $A$ with $A^2 = I_2$;
        \item the upper triangular $n \times n$ matrices whose diagonal entries are all equal to $1$;
        \item the $n \times n$ permutation matrices. How many of them are there?
    \end{enumerate}
\end{problem}

\begin{problem}\label{pr:inverse-triangular}
    Use the Gauss--Jordan method to find the inverse of $\begin{pmatrix} 1 & 2 & 3 \\ 0 & 1 & 4 \\ 0 & 0 & 1 \end{pmatrix}$. Then explain why the inverse of every invertible upper triangular matrix is again upper triangular.
\end{problem}

\begin{problem}\label{pr:orthogonal2}
    Find all orthogonal $2 \times 2$ matrices. (Hint: for $Q = \begin{pmatrix} a & b \\ c & d \end{pmatrix}$, the condition $Q^\top Q = I_2$ means $a^2 + c^2 = 1$, $b^2 + d^2 = 1$ and $ab + cd = 0$. Write $(a, c) = (\cos\theta, \sin\theta)$.) Which of them are rotations, and which are reflections?
\end{problem}

\subsection{Trace and Determinant}

\begin{problem}[$\star$]\label{pr:tridiagonal}
    Let $D_n$ be the determinant of the $n \times n$ matrix with $2$s on the diagonal, $1$s directly above and below the diagonal, and $0$s everywhere else; for example,
    \[D_3 = \det \begin{pmatrix} 2 & 1 & 0 \\ 1 & 2 & 1 \\ 0 & 1 & 2 \end{pmatrix}.\]
    \begin{enumerate}[label=(\alph*)]
        \item Compute $D_1$, $D_2$ and $D_3$.
        \item Expand along the first row to show that $D_n = 2D_{n-1} - D_{n-2}$ for $n \geq 3$.
        \item Find a formula for $D_n$.
    \end{enumerate}
\end{problem}

\begin{problem}[\mandatory]\label{pr:det4}
    Compute the determinant of
    \[A = \begin{pmatrix} 2 & 0 & 1 & 3 \\ 1 & 1 & 0 & 2 \\ 0 & 3 & 1 & 1 \\ 1 & 0 & 2 & 1 \end{pmatrix}\]
    in two ways: by cofactor expansion along a well-chosen row or column, and by row reduction. Is $A$ invertible?
\end{problem}

\begin{problem}[\mandatory]\label{pr:det-rules}
    Let $A$ and $B$ be $3 \times 3$ matrices with $\det(A) = 2$ and $\det(B) = -3$. Compute
    \[\det(AB), \qquad \det(A^\top B^2), \qquad \det(3A), \qquad \det(A^{-1}B), \qquad \det(BAB^{-1}).\]
    Let $C$ be obtained from $A$ by swapping rows $1$ and $3$ and then adding $5$ times row $1$ to row $2$. What is $\det(C)$?
\end{problem}

\begin{problem}\label{pr:vandermonde}
    \begin{enumerate}[label=(\alph*)]
        \item Use row operations to show that
        \[\det \begin{pmatrix} 1 & a & a^2 \\ 1 & b & b^2 \\ 1 & c & c^2 \end{pmatrix} = (b - a)(c - a)(c - b).\]
        \item Conclude that for distinct numbers $a$, $b$, $c$ and arbitrary values $y_1$, $y_2$, $y_3$, there is exactly one polynomial $p \in \mathbb{R}[x]_{\leq 2}$ with $p(a) = y_1$, $p(b) = y_2$ and $p(c) = y_3$. Compare with Problem~\ref{pr:interpolation}.
    \end{enumerate}
\end{problem}

\begin{problem}[$\star$]\label{pr:lambda}
    Let $A = \begin{pmatrix} 3 & 5 & 3 \\ 1 & 7 & 3 \\ 1 & 2 & 8 \end{pmatrix}$.
    \begin{enumerate}[label=(\alph*)]
        \item For which of the numbers $\lambda = 2, 3, 5$ is the matrix $A - \lambda I_3$ singular?
        \item For each such $\lambda$, find a nonzero vector $x$ with $Ax = \lambda x$.
        \item Without computing any determinant, find another number $\lambda$ and a nonzero vector $x$ with $Ax = \lambda x$. (Hint: add up the entries in each row of $A$.)
    \end{enumerate}
\end{problem}

\begin{problem}[$\star$]\label{pr:upper5}
    Let
    \[A = \begin{pmatrix} 1 & 2 & 3 & 4 & 5 \\ 0 & 2 & 3 & 4 & 5 \\ 0 & 0 & 3 & 4 & 5 \\ 0 & 0 & 0 & 4 & 5 \\ 0 & 0 & 0 & 0 & 5 \end{pmatrix}.\]
    \begin{enumerate}[label=(\alph*)]
        \item Compute $\det(A)$. Is $A$ invertible? Can $A$ be transformed into $I_5$ by elementary row operations?
        \item Find the last row of $A^2$ without computing all of $A^2$.
        \item Explain why $A - I_5$ and $A - 2I_5$ are singular. Find nonzero vectors $x$ and $y$ with $Ax = x$ and $Ay = 2y$.
    \end{enumerate}
\end{problem}

\begin{problem}\label{pr:area}
    \begin{enumerate}[label=(\alph*)]
        \item Find the area of the triangle with vertices $(1, 1)$, $(4, 2)$ and $(2, 5)$. (Hint: move one vertex to the origin.)
        \item Let $T(x, y) = (2x - y, x + 3y)$. Find the area of the image under $T$ of the unit square and of the unit disk.
        \item Show that the image of the unit disk under $T(x, y) = (ax, by)$, with $a, b > 0$, is the ellipse $\frac{x^2}{a^2} + \frac{y^2}{b^2} \leq 1$, and conclude that its area is $\pi ab$.
        \item Which of the transformations of Problem~\ref{pr:gallery} reverse the orientation? Compare with your sketches of the letter F.
    \end{enumerate}
\end{problem}

\begin{problem}\label{pr:fibonacci}
    The Fibonacci numbers are defined by $F_0 = 0$, $F_1 = 1$ and $F_{n+1} = F_n + F_{n-1}$. Let $M = \begin{pmatrix} 1 & 1 \\ 1 & 0 \end{pmatrix}$.
    \begin{enumerate}[label=(\alph*)]
        \item Show by induction that $M^n = \begin{pmatrix} F_{n+1} & F_n \\ F_n & F_{n-1} \end{pmatrix}$ for $n \geq 1$, and compute $M^{10}$.
        \item Use the product rule for determinants to prove Cassini's identity $F_{n+1}F_{n-1} - F_n^2 = (-1)^n$. Check it for $n = 10$.
    \end{enumerate}
\end{problem}

\begin{problem}\label{pr:trace}
    \begin{enumerate}[label=(\alph*)]
        \item Find $2 \times 2$ matrices $A$ and $B$ with $AB - BA = \begin{pmatrix} 1 & 0 \\ 0 & -1 \end{pmatrix}$. Why is there no such pair with $AB - BA = I_2$?
        \item Show that if $\tr(A^\top A) = 0$, then $A = 0$.
        \item Let $A$ be an $n \times n$ matrix with $\tr(AX) = 0$ for every $n \times n$ matrix $X$. Show that $A = 0$. (Hint: try $X = e_{ji}$, the matrix with a single $1$ in position $(j, i)$.)
    \end{enumerate}
\end{problem}

\begin{problem}[$\star$]\label{pr:01-matrices}
    \begin{enumerate}[label=(\alph*)]
        \item How many of the $16$ matrices $\begin{pmatrix} a & b \\ c & d \end{pmatrix}$ with all entries in $\{0, 1\}$ are invertible?
        \item Show that a square matrix $A$ with integer entries has an inverse with integer entries if and only if $\det(A) = \pm 1$.
    \end{enumerate}
\end{problem}

\section{Null Space, Column Space and Rank}

\begin{problem}[\mandatory]\label{pr:bases-B}
    Let $B = \begin{pmatrix} 1 & 2 & 0 & 3 \\ 2 & 4 & 1 & 8 \\ 1 & 2 & 1 & 5 \end{pmatrix}$. Find $\RREF(B)$, the rank of $B$, and bases of $\Null(B)$, $\col(B)$ and $\row(B)$. Check that your answer agrees with the rank--nullity theorem.
\end{problem}

\begin{problem}[\mandatory, $\star$]\label{pr:onto}
    \begin{enumerate}[label=(\alph*)]
        \item Let $T: \mathbb{R}^5 \to \mathbb{R}^3$ be a surjective linear transformation. What is the dimension of $\ker(T)$?
        \item Can a linear transformation $\mathbb{R}^3 \to \mathbb{R}^5$ be surjective? Can it be injective?
        \item Can a $4 \times 6$ matrix $A$ have $\Null(A) = \{0\}$?
        \item Let $A$ be a $5 \times 3$ matrix such that $Ax = 0$ only for $x = 0$. What is the rank of $A$?
    \end{enumerate}
\end{problem}

\begin{problem}[$\star$]\label{pr:intersect}
    \begin{enumerate}[label=(\alph*)]
        \item Let $V$ and $W$ be $4$-dimensional subspaces of $\mathbb{R}^7$. Show that $V \cap W \neq \{0\}$. (Hint: put a basis of $V$ and a basis of $W$ as the columns of a $7 \times 8$ matrix, and apply the rank--nullity theorem.)
        \item More generally, show that if $V$ and $W$ are subspaces of $\mathbb{R}^n$ with $\dim V + \dim W > n$, then $V \cap W \neq \{0\}$.
        \item Find two $3$-dimensional subspaces of $\mathbb{R}^6$ whose intersection is $\{0\}$.
    \end{enumerate}
\end{problem}

\begin{problem}\label{pr:rankAB}
    \begin{enumerate}[label=(\alph*)]
        \item Show that $\rank(AB) \leq \rank(A)$ and $\rank(AB) \leq \rank(B)$ whenever $AB$ is defined. (Hint: $\col(AB) \subseteq \col(A)$ and $\Null(B) \subseteq \Null(AB)$.)
        \item Conclude that if $A$ is $3 \times 2$ and $B$ is $2 \times 3$, then the $3 \times 3$ matrix $AB$ is never invertible.
    \end{enumerate}
\end{problem}

\begin{problem}\label{pr:ATA}
    Let $A$ be an $m \times n$ matrix.
    \begin{enumerate}[label=(\alph*)]
        \item Show that $\Null(A^\top A) = \Null(A)$. (Hint: if $A^\top A x = 0$, then $x^\top A^\top A x = 0$; what is $x^\top A^\top A x$ in terms of the entries of $Ax$?)
        \item Conclude that $\rank(A^\top A) = \rank(A)$.
        \item Show that if the columns of $A$ are linearly independent, then $A^\top A$ is invertible. (This is why the normal equations $(A^\top A)w = A^\top y$ of linear regression have a unique solution.)
    \end{enumerate}
\end{problem}

\section{Dual Vector Spaces}

\begin{problem}[\mandatory]\label{pr:functionals}
    Which of the following maps are linear functionals? For those on $\mathbb{R}^n$, write them as $x \mapsto ax$ with a row vector $a$.
    \begin{enumerate}[label=(\alph*)]
        \item $\mathbb{R}^3 \to \mathbb{R}$, $x \mapsto 2x_1 - x_3$;
        \item $\mathbb{R}^3 \to \mathbb{R}$, $x \mapsto x_1 x_2$;
        \item $\mathbb{R}^{2 \times 2} \to \mathbb{R}$, $A \mapsto \tr(A)$;
        \item $\mathbb{R}^{2 \times 2} \to \mathbb{R}$, $A \mapsto \det(A)$;
        \item $\mathbb{R}[x]_{\leq 3} \to \mathbb{R}$, $p \mapsto p(2) - p'(0)$;
        \item $\mathbb{R}[x]_{\leq 3} \to \mathbb{R}$, $p \mapsto p(0)^2$.
    \end{enumerate}
\end{problem}

\begin{problem}[\mandatory]\label{pr:dual-basis-r2}
    \begin{enumerate}[label=(\alph*)]
        \item Find the dual basis $b^1, b^2$ of the basis $b_1 = (2, 1)^\top$, $b_2 = (1, 1)^\top$ of $\mathbb{R}^2$, and check that $b^i(b_j) = \delta_{ij}$.
        \item Write the functional $\varphi(x) = 3x_1 + x_2$ as a linear combination of $b^1$ and $b^2$, without solving a system of equations.
    \end{enumerate}
\end{problem}

\begin{problem}[$\star$]\label{pr:lagrange-simpson}
    Let $V = \mathbb{R}[x]_{\leq 2}$ and consider the nodes $0$, $1$, $2$.
    \begin{enumerate}[label=(\alph*)]
        \item Find the Lagrange polynomials $L_0, L_1, L_2 \in V$ with $L_j(k) = \delta_{jk}$ for $j, k \in \{0, 1, 2\}$.
        \item Write $p(x) = x^2 + 1$ as a linear combination of $L_0, L_1, L_2$.
        \item The functional $I(p) = \int_0^2 p(x)\, dx$ on $V$ is a linear combination of the evaluation functionals $\operatorname{ev}_0$, $\operatorname{ev}_1$, $\operatorname{ev}_2$. Find the coefficients by computing $I(L_0)$, $I(L_1)$ and $I(L_2)$. (You have discovered Simpson's rule for numerical integration, which is exact for polynomials of degree at most $2$.)
    \end{enumerate}
\end{problem}

\begin{problem}[\mandatory]\label{pr:dual-map}
    Let $T: \mathbb{R}^3 \to \mathbb{R}^2$, $T(x, y, z) = (x + y, \ y - z)$, with matrix $A$.
    \begin{enumerate}[label=(\alph*)]
        \item For the functional $\psi$ on $\mathbb{R}^2$ given by the row $(2, -1)$, compute $T^*(\psi)$ as a row vector.
        \item Check that $T^*(\psi)(v) = \psi(T(v))$ for $v = (1, 2, 3)^\top$.
        \item Find the matrix of $T^*$ with respect to the duals of the standard bases, and compare it with $A$.
    \end{enumerate}
\end{problem}

\begin{problem}\label{pr:annihilator}
    \begin{enumerate}[label=(\alph*)]
        \item Find a basis of the annihilator of $U = \Span\{(1, 2, 0)^\top, (0, 1, 1)^\top\}$ in $(\mathbb{R}^3)^*$, and check that $\dim U + \dim U^\circ = 3$.
        \item Find a basis of the annihilator of $W = \Span\{(1, 0, 1, 0)^\top, (0, 1, 0, 1)^\top\}$ in $(\mathbb{R}^4)^*$.
    \end{enumerate}
\end{problem}

\begin{problem}[$\star$]\label{pr:certificate}
    \begin{enumerate}[label=(\alph*)]
        \item Show that the system $x_1 + 2x_2 = 1$, $2x_1 + 4x_2 = 3$ has no solution by finding a row $y$ with $yA = 0$ and $yb \neq 0$.
        \item For which $b = (b_1, b_2, b_3)^\top$ does the system with the matrix $A = \begin{pmatrix} 1 & 2 \\ 2 & 4 \\ 3 & 6 \end{pmatrix}$ have a solution? Answer by finding all rows $y$ with $yA = 0$ and using the Fredholm alternative.
    \end{enumerate}
\end{problem}

\begin{problem}\label{pr:double-dual}
    Let $V = \mathbb{R}^2$, and identify $V^*$ with row vectors. Describe the element $\iota(e_1) \in V^{**}$, i.e., compute $\iota(e_1)(\varphi)$ for the functional $\varphi$ given by a row $(a_1, a_2)$. Then show directly that $\iota(v) = 0$ implies $v = 0$.
\end{problem}

\section{Quotient Spaces and the Isomorphism Theorems}

\begin{problem}[\mandatory]\label{pr:cosets}
    Let $U = \Span\{(2, 1)^\top\} \subseteq \mathbb{R}^2$.
    \begin{enumerate}[label=(\alph*)]
        \item Decide whether $(3, 1)^\top + U = (1, 0)^\top + U$, whether $(5, 3)^\top + U = (1, 1)^\top + U$, and whether $(1, 2)^\top + U = U$.
        \item Show that the cosets of $U$ are the lines $x_1 - 2x_2 = c$, and find the unique representative of each coset on the $x_1$-axis.
        \item Find an explicit isomorphism $\mathbb{R}^2/U \to \mathbb{R}$.
    \end{enumerate}
\end{problem}

\begin{problem}\label{pr:quotient-dim}
    Find the dimensions of the following quotient spaces, and describe each of them as a familiar vector space:
    \begin{enumerate}[label=(\alph*)]
        \item $\mathbb{R}^5/U$, where $\dim U = 2$;
        \item $\mathbb{R}^{3 \times 3}/\mathrm{Sym}_3$;
        \item $\mathbb{R}[x]_{\leq 4}/U$, where $U = \{p : p(1) = p(2) = 0\}$. (Hint: consider $p \mapsto (p(1), p(2))$.)
    \end{enumerate}
\end{problem}

\begin{problem}[\mandatory]\label{pr:first-iso}
    Use the first isomorphism theorem to identify the following quotient spaces.
    \begin{enumerate}[label=(\alph*)]
        \item $\mathbb{R}^3/P$, where $P$ is the plane $x + y + z = 0$. (Use $T(x, y, z) = x + y + z$.)
        \item $\mathbb{R}[x]_{\leq 3}/U$, where $U = \{p : p(0) = p(1) = 0\}$. What is $\dim U$?
        \item $\mathbb{R}^{n \times n}/\mathrm{Sym}_n$. (Use $T(A) = A - A^\top$, and find its kernel and its image.)
    \end{enumerate}
\end{problem}

\begin{problem}[\mandatory, $\star$]\label{pr:grassmann}
    \begin{enumerate}[label=(\alph*)]
        \item Let $V$ and $W$ be $4$-dimensional subspaces of a $7$-dimensional vector space. Which values can $\dim(V \cap W)$ take? Give an example for the smallest possible value.
        \item Let $U = \{x \in \mathbb{R}^4 : x_1 + x_2 + x_3 + x_4 = 0\}$ and $W = \{x \in \mathbb{R}^4 : x_1 - x_2 + x_3 - x_4 = 0\}$. Use the dimension formula to find $\dim(U \cap W)$, and then find a basis of $U \cap W$.
    \end{enumerate}
\end{problem}

\begin{problem}\label{pr:third-iso}
    Let $V = \mathbb{R}^4$, $U = \Span\{e_4\}$ and $W = \Span\{e_3, e_4\}$. Describe $V/U$, $W/U$ and $V/W$, and check the third isomorphism theorem by comparing dimensions.
\end{problem}

\begin{problem}[\mandatory, $\star$]\label{pr:softmax}
    The softmax function is $\operatorname{softmax}(z)_i = \frac{e^{z_i}}{\sum_j e^{z_j}}$ for $z \in \mathbb{R}^n$.
    \begin{enumerate}[label=(\alph*)]
        \item Show that $\operatorname{softmax}(z + c\mathbf{1}) = \operatorname{softmax}(z)$ for every $c \in \mathbb{R}$, so that softmax is a function on $\mathbb{R}^n/\Span\{\mathbf{1}\}$.
        \item Compute $\operatorname{softmax}(1, 2, 3)$ and $\operatorname{softmax}(1001, 1002, 1003)$. Why can a computer evaluate the second one only after subtracting a constant from all entries?
        \item Is softmax injective on $\mathbb{R}^n/\Span\{\mathbf{1}\}$? (Hint: if $\operatorname{softmax}(z) = \operatorname{softmax}(w)$, compare $z_i - z_j$ and $w_i - w_j$.)
    \end{enumerate}
\end{problem}

\begin{problem}\label{pr:quotient-dual}
    Let $U = \Span\{(1, 1, 1)^\top\} \subseteq \mathbb{R}^3$. Describe all linear functionals on $\mathbb{R}^3/U$ as rows $a = (a_1, a_2, a_3)$, and check that they form a $2$-dimensional space, in agreement with $(V/U)^* \cong U^\circ$.
\end{problem}

\section{Eigenvalues and Eigenvectors}

\begin{problem}[\mandatory]\label{pr:eig2}
    Let $A = \begin{pmatrix} 1 & 2 \\ 3 & 2 \end{pmatrix}$.
    \begin{enumerate}[label=(\alph*)]
        \item Compute the characteristic polynomial of $A$, its eigenvalues, and its eigenspaces.
        \item Check that the eigenvalues add up to $\tr(A)$ and multiply to $\det(A)$.
        \item Diagonalize $A$, i.e., find $P$ and $D$ with $A = PDP^{-1}$.
        \item Find a formula for $A^k$, and check it for $k = 1$.
    \end{enumerate}
\end{problem}

\begin{problem}\label{pr:eig-triangular}
    Find the eigenvalues and eigenspaces of $L = \begin{pmatrix} 2 & 0 & 0 \\ 1 & 3 & 0 \\ 4 & -1 & 1 \end{pmatrix}$. Is $L$ diagonalizable?
\end{problem}

\begin{problem}[\mandatory]\label{pr:eig-defective}
    Let $N = \begin{pmatrix} 2 & 1 & 0 \\ 0 & 2 & 0 \\ 0 & 0 & 3 \end{pmatrix}$. Find the algebraic and geometric multiplicity of each eigenvalue of $N$, and decide whether $N$ is diagonalizable.
\end{problem}

\begin{problem}\label{pr:eig-geometric}
    Without computing any characteristic polynomial, use geometric intuition to find the eigenvalues and eigenspaces of the following linear transformations of the plane. Then check your answers with a computation.
    \begin{enumerate}[label=(\alph*)]
        \item the reflection in the $x$-axis;
        \item the orthogonal projection onto the line $y = x$, with matrix $\frac{1}{2}\begin{pmatrix} 1 & 1 \\ 1 & 1 \end{pmatrix}$;
        \item the shear $T(x, y) = (x + 2y, y)$;
        \item the rotation by $90^\circ$;
        \item $T(x, y) = (2x, 3y)$;
        \item $T(x, y) = (-x, -y)$.
    \end{enumerate}
    Which of these transformations are diagonalizable?
\end{problem}

\begin{problem}[$\star$]\label{pr:eig-gre}
    A real $3 \times 3$ matrix $A$ has the eigenvalues $-1$, $2$ and $4$.
    \begin{enumerate}[label=(\alph*)]
        \item Is $A$ diagonalizable? Is $A$ invertible?
        \item Compute $\tr(A)$, $\det(A)$, $\tr(A^2)$ and $\det(A^3)$.
        \item Find the eigenvalues of $A^{-1}$, $A^2$ and $A + 2I_3$.
    \end{enumerate}
\end{problem}

\begin{problem}[$\star$]\label{pr:eig-idempotent}
    Let $P$ be an $n \times n$ matrix with $P^2 = P$ (compare Problem~\ref{pr:idempotent}).
    \begin{enumerate}[label=(\alph*)]
        \item Show that every eigenvalue of $P$ is $0$ or $1$.
        \item Show that every nonzero vector of $\im(T_P)$ is an eigenvector for the eigenvalue $1$, and that every nonzero vector of $\ker(T_P)$ is an eigenvector for the eigenvalue $0$. Use Problem~\ref{pr:idempotent}(d) to conclude that $P$ is diagonalizable.
        \item Which of the following statements must be true: $P$ is invertible; $P$ is the identity or the zero matrix; $P$ is diagonalizable?
    \end{enumerate}
\end{problem}

\begin{problem}\label{pr:eig-props}
    Let $A$ and $B$ be $n \times n$ matrices.
    \begin{enumerate}[label=(\alph*)]
        \item Show that $A$ and $A^\top$ have the same eigenvalues, and give an example where they have different eigenvectors.
        \item Show that if $A$ is invertible, then $AB$ and $BA$ are similar, and hence have the same eigenvalues. (Hint: $BA = A^{-1}(AB)A$.)
        \item Show that every real eigenvalue of an orthogonal matrix $Q$ is $1$ or $-1$. Hint: if $Qv = \lambda v$, compute $(Qv)^\top (Qv)$ in two ways.
        \item Show that if $A^k = 0$ for some $k \geq 1$, then $0$ is the only eigenvalue of $A$.
        \item Show that $\begin{pmatrix} 1 & 1 \\ 0 & 1 \end{pmatrix}$ and $I_2$ have the same characteristic polynomial, but are not similar.
    \end{enumerate}
\end{problem}

\begin{problem}[$\star$]\label{pr:eig-recurrence}
    The sequence $a_0 = 0$, $a_1 = 1$, $a_{n+1} = a_n + 2a_{n-1}$ starts with $0, 1, 1, 3, 5, 11, \ldots$.
    \begin{enumerate}[label=(\alph*)]
        \item Find a $2 \times 2$ matrix $M$ with $(a_{n+1}, a_n)^\top = M (a_n, a_{n-1})^\top$.
        \item Find the eigenvalues and eigenvectors of $M$, and derive a formula for $a_n$, as for Binet's formula in the lecture.
        \item Compute $a_{10}$ with your formula.
    \end{enumerate}
\end{problem}

\begin{problem}[\mandatory]\label{pr:eig-markov}
    Every day, a visitor of a website either only browses or makes a purchase. A browsing visitor browses again the next day with probability $0.7$, and a buying visitor buys again with probability $0.6$. With $x_k = (b_k, p_k)^\top$ the probabilities of browsing and buying on day $k$, we have $x_{k+1} = Ax_k$ with
    \[A = \begin{pmatrix} 0.7 & 0.4 \\ 0.3 & 0.6 \end{pmatrix}.\]
    \begin{enumerate}[label=(\alph*)]
        \item Find the eigenvalues and eigenvectors of $A$, and the stationary distribution.
        \item For $x_0 = (1, 0)^\top$, find a formula for $x_k$, and check it for $k = 2$. What happens as $k \to \infty$?
    \end{enumerate}
\end{problem}

\begin{problem}\label{pr:eig-power}
    Let $A = \begin{pmatrix} 2 & 1 \\ 1 & 2 \end{pmatrix}$ and $x_0 = (1, 0)^\top$.
    \begin{enumerate}[label=(\alph*)]
        \item Compute $Ax_0$, $A^2x_0$, $A^3x_0$ and $A^4x_0$. What happens to the ratio of the two entries?
        \item Explain your observation by writing $x_0$ as a linear combination of eigenvectors of $A$, and find a formula for $A^k x_0$.
    \end{enumerate}
\end{problem}

\begin{problem}[$\star$]\label{pr:eig-complex}
    \begin{enumerate}[label=(\alph*)]
        \item Show that the matrix $M(a + bi) = \begin{pmatrix} a & -b \\ b & a \end{pmatrix}$ of Problem~\ref{pr:complex-matrix} has the complex eigenvalues $a \pm bi$.
        \item Find the complex eigenvalues of $\begin{pmatrix} 3 & -4 \\ 4 & 3 \end{pmatrix}$, and write this matrix as $r R_\theta$ with $r > 0$ and a rotation $R_\theta$. How are $r$ and $\theta$ related to the eigenvalues?
    \end{enumerate}
\end{problem}

\end{document}
